\chapter{Creating Equations}\label{Chap1_Equations}

In this chapter, we show the different ways of writing and citing equations. Here is the first equation:
\color{black}
\begin{equation}
	\label{Newton2ndLaw}
	F = ma \,. 
\end{equation} 

And here is how we remove the equation number

\begin{equation*}
	\label{Newton2ndLaw}
	F = ma \,. 
\end{equation*} 



\color{blue}
Citing an equation follows the syntax below:

``According to Equation (\ref{Newton2ndLaw}) ...''

\noindent Use Equation with a capital `E' even if you are citing in the middle of the sentence. The equation number should be in parentheses.
\color{red}
 Here are some \textit{wrong} syntax that you should avoid using: 

	\begin{singlespace}
\color{red}
\begin{enumerate}
	\item According to equation \ref{Newton2ndLaw}
	\item According to eq \ref{Newton2ndLaw}
	\item According to Eq \ref{Newton2ndLaw}
	\item According to Eq. \ref{Newton2ndLaw}
	\item According to equation (\ref{Newton2ndLaw})
	\item According to Eq. (\ref{Newton2ndLaw})
\end{enumerate}
 \color{black}

\end{singlespace}


\color{black}
If you want to divide your equation into two sub-equations, then you can use the following commands: 
\begin{subequations}
	\begin{equation}
		\label{Body1}
	F_{1} = m_{1}a_{1} \,. 
	\end{equation}
	\begin{equation}
		\label{Body2}
	F_{2} = m_{2}a_{2} \,.  
	\end{equation}
\end{subequations}



You can now cite these equations as following: 

\color{blue} ``According to Equations (\ref{Body1}) and (\ref{Body2}) ...''\color{black}
